% codex_notes.tex — 工作區共用的講義版型 v2（2026-10-09 從量導講義 v2 抽出）
% 用於：量導 lectures/template/quantum_notes_v2.sty、熱物理 lectures/template/thermal_notes_v2.sty。
% （群論 kuo_notes_v2.sty 是數學講義的變體，差異多，仍獨立。）
% 各專案的 .sty 先設參數，再用下面這行載入本檔，之後只放記號巨集與專案特有的區塊：
%   \IfFileExists{codex_notes.tex}{% codex_notes.tex — 工作區共用的講義版型 v2（2026-10-09 從量導講義 v2 抽出）
% 用於：量導 lectures/template/quantum_notes_v2.sty、熱物理 lectures/template/thermal_notes_v2.sty。
% （群論 kuo_notes_v2.sty 是數學講義的變體，差異多，仍獨立。）
% 各專案的 .sty 先設參數，再用下面這行載入本檔，之後只放記號巨集與專案特有的區塊：
%   \IfFileExists{codex_notes.tex}{% codex_notes.tex — 工作區共用的講義版型 v2（2026-10-09 從量導講義 v2 抽出）
% 用於：量導 lectures/template/quantum_notes_v2.sty、熱物理 lectures/template/thermal_notes_v2.sty。
% （群論 kuo_notes_v2.sty 是數學講義的變體，差異多，仍獨立。）
% 各專案的 .sty 先設參數，再用下面這行載入本檔，之後只放記號巨集與專案特有的區塊：
%   \IfFileExists{codex_notes.tex}{% codex_notes.tex — 工作區共用的講義版型 v2（2026-10-09 從量導講義 v2 抽出）
% 用於：量導 lectures/template/quantum_notes_v2.sty、熱物理 lectures/template/thermal_notes_v2.sty。
% （群論 kuo_notes_v2.sty 是數學講義的變體，差異多，仍獨立。）
% 各專案的 .sty 先設參數，再用下面這行載入本檔，之後只放記號巨集與專案特有的區塊：
%   \IfFileExists{codex_notes.tex}{\input{codex_notes.tex}}{\input{"D:/Codex workplace/_templates/styles/codex_notes.tex"}}
% 參數（都可省略）：\codexsecsep（節號分隔，預設 -）、\codexlecturetag{章號}（頁首右側，預設 L03）、\codextoclabelwidth（預設 2.6em）。
% 版面：A4、10pt（主檔 \documentclass[10pt,a4paper]{article}）、四邊 2cm、行距 1.3、Noto Sans TC、段落不縮排；深藍標題、藍色節號、淺橘重點結果。
% 巨集：\setuplecture{章號}{章名}、\lecturehead{小字}{章名}、\booksection{中文}{English}、\plainsection、\chaptergoal、\keyresult（只放數學式）、
% definition／theorem／note／derivation／example／solution／extension[標題]、\coursekeepblocks；記號 \dd、\e、\avg、\Order。
\RequirePackage[no-math]{fontspec}
\RequirePackage{xeCJK}
% 英文撇號：' 經 TeX 連字變成 U+2019，xeCJK 預設把 ‘’ 當全形標點，Noto Sans TC 的這兩個字形也是全寬，
% Ampère's 的撇號後面因此多一段空隙。把 ‘’ 改回西文字元，主字型加 pwid 取比例寬字形（只作用於西文字型，
% 中文標點由 CJK 字型排，不受影響）。中文引號用「」，“” 維持原狀。2026-10-22
\defaultfontfeatures+[\rmfamily]{RawFeature=+pwid}
\xeCJKDeclareCharClass{Default}{"2018,"2019}
% 字型資料夾：載入本檔前 \def\codexfontpath{<資料夾>/} 可改；設成空白則依檔名在 TeX 搜尋路徑與目前資料夾找。2026-11-15
\providecommand{\codexfontpath}{D:/Codex workplace/Noto_Sans_TC/static/}
\setmainfont{NotoSansTC-Regular.ttf}[
  Path=\codexfontpath,
  BoldFont=NotoSansTC-Bold.ttf,
  ItalicFont=NotoSansTC-Regular.ttf,
  BoldItalicFont=NotoSansTC-Bold.ttf
]
\setCJKmainfont{NotoSansTC-Regular.ttf}[
  Path=\codexfontpath,
  BoldFont=NotoSansTC-Bold.ttf,
  ItalicFont=NotoSansTC-Regular.ttf,
  BoldItalicFont=NotoSansTC-Bold.ttf
]
\setmonofont{NotoSansTC-Regular.ttf}[
  Path=\codexfontpath,
  BoldFont=NotoSansTC-Bold.ttf
]
\setCJKmonofont{NotoSansTC-Regular.ttf}[
  Path=\codexfontpath,
  BoldFont=NotoSansTC-Bold.ttf
]
\RequirePackage{amsmath,amssymb,amsthm,mathtools}
\RequirePackage[a4paper,margin=2cm]{geometry}
\RequirePackage{booktabs,longtable,array,tabularx}
\RequirePackage{graphicx,float}
\RequirePackage{enumitem,setspace,xcolor,fancyhdr,needspace,etoolbox}
\RequirePackage[most]{tcolorbox}
\RequirePackage{titlesec,titletoc}
\RequirePackage{caption}
\RequirePackage{tikz,pgfplots}
\RequirePackage[unicode]{hyperref}
\pgfplotsset{compat=1.18}
\usetikzlibrary{arrows.meta,calc,positioning,decorations.pathreplacing}
\usepgfplotslibrary{fillbetween}

% ---- 色票（與共用詳解版型相同，另加綠、紫各一）----
\definecolor{qblue}{HTML}{1D5FA8}
\definecolor{qnavy}{HTML}{173D6B}
\definecolor{qorange}{HTML}{EB6834}
\definecolor{qorangebg}{HTML}{FFF4E8}
\definecolor{qgraybg}{HTML}{F3F4F6}
\definecolor{qrule}{HTML}{D8DDE3}
\definecolor{qmuted}{HTML}{6B7280}
\definecolor{qtext}{HTML}{1F2329}
\definecolor{qgreen}{HTML}{2E7D4F}
\definecolor{qpurple}{HTML}{6B3FA0}
% 舊版型用到的灰色名稱（圖檔可能引用）
\colorlet{qgray}{qmuted}
\colorlet{qlight}{qgraybg}
\hypersetup{colorlinks=true,linkcolor=qnavy,urlcolor=qblue,citecolor=qblue}

% ---- 段落 ----
\setstretch{1.3}
\setlength{\parindent}{0pt}
\setlength{\parskip}{0.55em}
\setlength{\headheight}{16pt}
\setlength{\headsep}{0.45cm}
\setlength{\footskip}{1cm}
\setlist{itemsep=0.2em,topsep=0.3em}
\setlist[itemize]{label=\textcolor{qblue}{\small\textbullet}}
\setlist[enumerate]{label=\textcolor{qblue}{\arabic*.}}
\renewcommand{\arraystretch}{1.16}
\allowdisplaybreaks[2]
\raggedbottom
\widowpenalty=10000
\clubpenalty=10000
\emergencystretch=3em
\AtBeginDocument{\color{qtext}}

% ---- 圖說 ----
\renewcommand{\figurename}{圖}
\renewcommand{\tablename}{表}
\captionsetup{font={small,color=qmuted},labelfont={bf,color=qnavy},labelsep=quad,
  justification=justified,singlelinecheck=true,skip=6pt}

% ---- 章設定：節號 3-1（分隔符 \codexsecsep）、式號與區塊號 (3.k)；頁首右側 \codexlecturetag{章號}　章名 ----
\providecommand{\codexsecsep}{-}
\providecommand{\codexlecturetag}[1]{L\ifnum#1<10 0\fi#1}
\providecommand{\codextoclabelwidth}{2.6em}
\newcommand{\lecturechapter}{0}
\newcommand{\lecturetitle}{}
\newcommand{\setuplecture}[2]{%
  \renewcommand{\lecturechapter}{#1}%
  \renewcommand{\lecturetitle}{#2}%
  \renewcommand{\thesection}{#1\codexsecsep\arabic{section}}%
  \renewcommand{\theequation}{#1.\arabic{equation}}%
  \renewcommand{\thefigure}{#1.\arabic{figure}}%
  \renewcommand{\thetable}{#1.\arabic{table}}%
  \fancyhead[L]{\color{qmuted}\small\coursename}%
  \fancyhead[R]{\color{qmuted}\small \codexlecturetag{#1}\quad #2}%
}

% ---- 頁首頁尾 ----
\pagestyle{fancy}
\fancyhf{}
\fancyfoot[C]{\color{qmuted}\small\thepage}
\renewcommand{\headrulewidth}{0pt}

% ---- 標題 ----
\titleformat{\section}[block]{\color{qnavy}\Large\bfseries\boldmath}
  {\textcolor{qblue}{\thesection}}{0.8em}{}
\titleformat{name=\section,numberless}[block]{\color{qnavy}\Large\bfseries\boldmath}{}{0pt}{}
\titlespacing*{\section}{0pt}{2.2em}{0.8em}
\titleformat{\subsection}{\color{qnavy}\normalsize\bfseries\boldmath}{}{0pt}{}
\titlespacing*{\subsection}{0pt}{1.2em}{0.25em}

% 課本節：\booksection{中文標題}{English title}；英文副標放同一行，放不下整段換行，不斷字
\newsavebox{\qsubtitlebox}
\DeclareRobustCommand{\qbooksubtitle}[1]{%
  \sbox{\qsubtitlebox}{\normalfont\small(#1)}%
  \ifdim\wd\qsubtitlebox<\dimexpr\linewidth-4em\relax
    \hskip 0pt plus 1fil\penalty0\hskip 0.8em plus -1fil\relax
    {\normalfont\small\color{qmuted}#1}%
  \else
    \enspace{\normalfont\small\color{qmuted}
      \ifcsname l@nohyphenation\endcsname\language\csname l@nohyphenation\endcsname\fi
      #1}%
  \fi}
\newcommand{\booksection}[2]{\section[#1]{#1\qbooksubtitle{#2}}}
\newcommand{\plainsection}[1]{\section*{#1}\addcontentsline{toc}{section}{#1}}

% ---- 目錄 ----
\renewcommand{\contentsname}{目錄}
\setcounter{tocdepth}{1}
\titlecontents{section}[\codextoclabelwidth]{\addvspace{0.35em}\color{qtext}}
  {\contentslabel[\textcolor{qblue}{\bfseries\thecontentslabel}]{\codextoclabelwidth}}{}
  {\color{qmuted}\titlerule*[0.6pc]{.}\contentspage}

% ---- 章首標頭（取代 \maketitle）----
\newcommand{\lecturehead}[2]{%
  {\color{qmuted}\small #1\par}\vspace{2pt}
  {\color{qnavy}\fontsize{18pt}{24pt}\selectfont\bfseries #2\par}\vspace{6pt}
  {\color{qrule}\rule{\linewidth}{0.6pt}\par}\vspace{0.6em}}

% ---- 記號巨集（共用的；其餘在專案的 .sty）----
\newcommand{\dd}{\mathop{}\!\mathrm{d}}
\newcommand{\e}{\mathrm{e}}
\newcommand{\avg}[1]{\left\langle #1\right\rangle}
\newcommand{\Order}{\mathcal O}

% ---- 區塊共用計數器：定義、定理、推導、例題共用 (3.k) ----
\newcounter{qblock}
\renewcommand{\theqblock}{\lecturechapter.\arabic{qblock}}
\newcommand{\qtitle}[1]{\ifstrempty{#1}{}{\enspace#1}}

% ---- 不讓標題落在頁底 ----
\makeatletter
\newcommand{\qblockspace}[1]{\if@nobreak\else\Needspace*{#1}\fi}
\makeatother
\newcommand{\coursekeepblocks}{%
  \pretocmd{\section}{\qblockspace{12\baselineskip}}{}{}%
  \pretocmd{\subsection}{\qblockspace{8\baselineskip}}{}{}%
}

% ---- 區塊 ----
% 框內 \label 記到區塊編號：tcolorbox 在 \sbox 裡做 \refstepcounter，\@currentlabel 出了那個群組就還原，
% 框內的 \label 因此記成外層的小節號（DE L02 遇過，2026-10-21）。\qsetlabel 在框的正文開頭重設一次；
% 自己寫 before upper 的編號框要把它放進去。沒有編號的框不受影響。
\makeatletter
\newcommand{\qsetlabel}{\ifdefined\tcbcounter\ifx\tcbcounter\@empty\else
  \protected@edef\@currentlabel{\lecturechapter.\thetcbcounter}\fi\fi}
\makeatother
\tcbset{qbase/.style={enhanced,breakable,arc=2pt,boxrule=0.6pt,
  left=8pt,right=8pt,top=5pt,bottom=5pt,before skip=0.9em,after skip=0.9em,
  fonttitle=\bfseries\small,toptitle=2pt,bottomtitle=2pt,
  before upper={\setlength{\parskip}{0.45em}\qsetlabel}}}
\tcbset{qbar/.style={enhanced,breakable,frame hidden,sharp corners,
  left=10pt,right=8pt,top=5pt,bottom=5pt,before skip=0.8em,after skip=0.8em,
  before upper={\setlength{\parskip}{0.45em}\qsetlabel}}}
\newtcolorbox[use counter=qblock]{definition}[1][]{qbase,
  colframe=qblue,colback=qblue!4,coltitle=white,colbacktitle=qblue,
  title={定義 \lecturechapter.\thetcbcounter\qtitle{#1}}}
\newtcolorbox[use counter=qblock]{theorem}[1][]{qbase,
  colframe=qpurple,colback=qpurple!4,coltitle=white,colbacktitle=qpurple,
  title={定理 \lecturechapter.\thetcbcounter\qtitle{#1}}}
\newtcolorbox{note}[1][]{qbase,
  colframe=qorange,colback=qorangebg,coltitle=white,colbacktitle=qorange,
  title={注意\qtitle{#1}}}
\newtcolorbox[use counter=qblock]{derivation}[1][]{qbar,colback=white,top=1pt,bottom=1pt,
  borderline west={1.2pt}{0pt}{qrule},
  before upper={\setlength{\parskip}{0.45em}\qsetlabel{\small\bfseries\color{qnavy}推導 \lecturechapter.\thetcbcounter\qtitle{#1}}\par}}
\newtcolorbox[use counter=qblock]{example}[1][]{qbar,colback=white,top=1pt,bottom=1pt,
  borderline west={1.2pt}{0pt}{qgreen!60},
  before upper={\setlength{\parskip}{0.45em}\qsetlabel{\small\bfseries\color{qgreen}例 \lecturechapter.\thetcbcounter\qtitle{#1}}\par}}
\newtcolorbox{extension}[1][]{qbar,colback=qgraybg,fontupper=\small\color{qmuted},
  before upper={{\bfseries 延伸\qtitle{#1}}\par}}
% 重點結果：與詳解相同的淺橘底＋左側橘條
\newtcolorbox{qresult}{enhanced,frame hidden,sharp corners,colback=qorangebg,
  borderline west={2.2pt}{0pt}{qorange},left=10pt,right=10pt,top=4pt,bottom=4pt,
  before skip=0.7em,after skip=0.7em,halign=center}
\newcommand{\keyresult}[1]{\par\begin{qresult}\makebox[\linewidth]{$\displaystyle #1$}\end{qresult}}
% 章首問題：淺灰底
\newcommand{\chaptergoal}[1]{%
  \begin{tcolorbox}[qbar,colback=qgraybg]{\small\bfseries\color{qnavy}本章問題}\par #1\end{tcolorbox}}

% 例題的解：接在 example 之後，綠色小標，不加框
\newenvironment{solution}
  {\par\smallskip\noindent{\small\bfseries\color{qgreen}解}\enspace\ignorespaces}
  {\par\medskip}
% 例題、推導前留空間，避免小標落在頁底
\BeforeBeginEnvironment{example}{\qblockspace{6\baselineskip}}
\BeforeBeginEnvironment{derivation}{\qblockspace{6\baselineskip}}
\BeforeBeginEnvironment{definition}{\qblockspace{5\baselineskip}}
\BeforeBeginEnvironment{theorem}{\qblockspace{5\baselineskip}}
% 可跨頁的框剛好在頁底結束時，xdvipdfmx 的顏色堆疊可能錯位，之後的正文變成白字
% （Neamen L05_v1 遇過，2026-10-15）。框的環境群組結束後重設一次正文顏色；
% 放在框的 after 鍵裡沒有用，因為群組結束時顏色又被恢復成錯位的狀態。
\AfterEndEnvironment{definition}{\normalcolor}
\AfterEndEnvironment{theorem}{\normalcolor}
\AfterEndEnvironment{property}{\normalcolor}
\AfterEndEnvironment{note}{\normalcolor}
\AfterEndEnvironment{derivation}{\normalcolor}
\AfterEndEnvironment{example}{\normalcolor}
\AfterEndEnvironment{extension}{\normalcolor}
}{\input{"D:/Codex workplace/_templates/styles/codex_notes.tex"}}
% 參數（都可省略）：\codexsecsep（節號分隔，預設 -）、\codexlecturetag{章號}（頁首右側，預設 L03）、\codextoclabelwidth（預設 2.6em）。
% 版面：A4、10pt（主檔 \documentclass[10pt,a4paper]{article}）、四邊 2cm、行距 1.3、Noto Sans TC、段落不縮排；深藍標題、藍色節號、淺橘重點結果。
% 巨集：\setuplecture{章號}{章名}、\lecturehead{小字}{章名}、\booksection{中文}{English}、\plainsection、\chaptergoal、\keyresult（只放數學式）、
% definition／theorem／note／derivation／example／solution／extension[標題]、\coursekeepblocks；記號 \dd、\e、\avg、\Order。
\RequirePackage[no-math]{fontspec}
\RequirePackage{xeCJK}
% 英文撇號：' 經 TeX 連字變成 U+2019，xeCJK 預設把 ‘’ 當全形標點，Noto Sans TC 的這兩個字形也是全寬，
% Ampère's 的撇號後面因此多一段空隙。把 ‘’ 改回西文字元，主字型加 pwid 取比例寬字形（只作用於西文字型，
% 中文標點由 CJK 字型排，不受影響）。中文引號用「」，“” 維持原狀。2026-10-22
\defaultfontfeatures+[\rmfamily]{RawFeature=+pwid}
\xeCJKDeclareCharClass{Default}{"2018,"2019}
% 字型資料夾：載入本檔前 \def\codexfontpath{<資料夾>/} 可改；設成空白則依檔名在 TeX 搜尋路徑與目前資料夾找。2026-11-15
\providecommand{\codexfontpath}{D:/Codex workplace/Noto_Sans_TC/static/}
\setmainfont{NotoSansTC-Regular.ttf}[
  Path=\codexfontpath,
  BoldFont=NotoSansTC-Bold.ttf,
  ItalicFont=NotoSansTC-Regular.ttf,
  BoldItalicFont=NotoSansTC-Bold.ttf
]
\setCJKmainfont{NotoSansTC-Regular.ttf}[
  Path=\codexfontpath,
  BoldFont=NotoSansTC-Bold.ttf,
  ItalicFont=NotoSansTC-Regular.ttf,
  BoldItalicFont=NotoSansTC-Bold.ttf
]
\setmonofont{NotoSansTC-Regular.ttf}[
  Path=\codexfontpath,
  BoldFont=NotoSansTC-Bold.ttf
]
\setCJKmonofont{NotoSansTC-Regular.ttf}[
  Path=\codexfontpath,
  BoldFont=NotoSansTC-Bold.ttf
]
\RequirePackage{amsmath,amssymb,amsthm,mathtools}
\RequirePackage[a4paper,margin=2cm]{geometry}
\RequirePackage{booktabs,longtable,array,tabularx}
\RequirePackage{graphicx,float}
\RequirePackage{enumitem,setspace,xcolor,fancyhdr,needspace,etoolbox}
\RequirePackage[most]{tcolorbox}
\RequirePackage{titlesec,titletoc}
\RequirePackage{caption}
\RequirePackage{tikz,pgfplots}
\RequirePackage[unicode]{hyperref}
\pgfplotsset{compat=1.18}
\usetikzlibrary{arrows.meta,calc,positioning,decorations.pathreplacing}
\usepgfplotslibrary{fillbetween}

% ---- 色票（與共用詳解版型相同，另加綠、紫各一）----
\definecolor{qblue}{HTML}{1D5FA8}
\definecolor{qnavy}{HTML}{173D6B}
\definecolor{qorange}{HTML}{EB6834}
\definecolor{qorangebg}{HTML}{FFF4E8}
\definecolor{qgraybg}{HTML}{F3F4F6}
\definecolor{qrule}{HTML}{D8DDE3}
\definecolor{qmuted}{HTML}{6B7280}
\definecolor{qtext}{HTML}{1F2329}
\definecolor{qgreen}{HTML}{2E7D4F}
\definecolor{qpurple}{HTML}{6B3FA0}
% 舊版型用到的灰色名稱（圖檔可能引用）
\colorlet{qgray}{qmuted}
\colorlet{qlight}{qgraybg}
\hypersetup{colorlinks=true,linkcolor=qnavy,urlcolor=qblue,citecolor=qblue}

% ---- 段落 ----
\setstretch{1.3}
\setlength{\parindent}{0pt}
\setlength{\parskip}{0.55em}
\setlength{\headheight}{16pt}
\setlength{\headsep}{0.45cm}
\setlength{\footskip}{1cm}
\setlist{itemsep=0.2em,topsep=0.3em}
\setlist[itemize]{label=\textcolor{qblue}{\small\textbullet}}
\setlist[enumerate]{label=\textcolor{qblue}{\arabic*.}}
\renewcommand{\arraystretch}{1.16}
\allowdisplaybreaks[2]
\raggedbottom
\widowpenalty=10000
\clubpenalty=10000
\emergencystretch=3em
\AtBeginDocument{\color{qtext}}

% ---- 圖說 ----
\renewcommand{\figurename}{圖}
\renewcommand{\tablename}{表}
\captionsetup{font={small,color=qmuted},labelfont={bf,color=qnavy},labelsep=quad,
  justification=justified,singlelinecheck=true,skip=6pt}

% ---- 章設定：節號 3-1（分隔符 \codexsecsep）、式號與區塊號 (3.k)；頁首右側 \codexlecturetag{章號}　章名 ----
\providecommand{\codexsecsep}{-}
\providecommand{\codexlecturetag}[1]{L\ifnum#1<10 0\fi#1}
\providecommand{\codextoclabelwidth}{2.6em}
\newcommand{\lecturechapter}{0}
\newcommand{\lecturetitle}{}
\newcommand{\setuplecture}[2]{%
  \renewcommand{\lecturechapter}{#1}%
  \renewcommand{\lecturetitle}{#2}%
  \renewcommand{\thesection}{#1\codexsecsep\arabic{section}}%
  \renewcommand{\theequation}{#1.\arabic{equation}}%
  \renewcommand{\thefigure}{#1.\arabic{figure}}%
  \renewcommand{\thetable}{#1.\arabic{table}}%
  \fancyhead[L]{\color{qmuted}\small\coursename}%
  \fancyhead[R]{\color{qmuted}\small \codexlecturetag{#1}\quad #2}%
}

% ---- 頁首頁尾 ----
\pagestyle{fancy}
\fancyhf{}
\fancyfoot[C]{\color{qmuted}\small\thepage}
\renewcommand{\headrulewidth}{0pt}

% ---- 標題 ----
\titleformat{\section}[block]{\color{qnavy}\Large\bfseries\boldmath}
  {\textcolor{qblue}{\thesection}}{0.8em}{}
\titleformat{name=\section,numberless}[block]{\color{qnavy}\Large\bfseries\boldmath}{}{0pt}{}
\titlespacing*{\section}{0pt}{2.2em}{0.8em}
\titleformat{\subsection}{\color{qnavy}\normalsize\bfseries\boldmath}{}{0pt}{}
\titlespacing*{\subsection}{0pt}{1.2em}{0.25em}

% 課本節：\booksection{中文標題}{English title}；英文副標放同一行，放不下整段換行，不斷字
\newsavebox{\qsubtitlebox}
\DeclareRobustCommand{\qbooksubtitle}[1]{%
  \sbox{\qsubtitlebox}{\normalfont\small(#1)}%
  \ifdim\wd\qsubtitlebox<\dimexpr\linewidth-4em\relax
    \hskip 0pt plus 1fil\penalty0\hskip 0.8em plus -1fil\relax
    {\normalfont\small\color{qmuted}#1}%
  \else
    \enspace{\normalfont\small\color{qmuted}
      \ifcsname l@nohyphenation\endcsname\language\csname l@nohyphenation\endcsname\fi
      #1}%
  \fi}
\newcommand{\booksection}[2]{\section[#1]{#1\qbooksubtitle{#2}}}
\newcommand{\plainsection}[1]{\section*{#1}\addcontentsline{toc}{section}{#1}}

% ---- 目錄 ----
\renewcommand{\contentsname}{目錄}
\setcounter{tocdepth}{1}
\titlecontents{section}[\codextoclabelwidth]{\addvspace{0.35em}\color{qtext}}
  {\contentslabel[\textcolor{qblue}{\bfseries\thecontentslabel}]{\codextoclabelwidth}}{}
  {\color{qmuted}\titlerule*[0.6pc]{.}\contentspage}

% ---- 章首標頭（取代 \maketitle）----
\newcommand{\lecturehead}[2]{%
  {\color{qmuted}\small #1\par}\vspace{2pt}
  {\color{qnavy}\fontsize{18pt}{24pt}\selectfont\bfseries #2\par}\vspace{6pt}
  {\color{qrule}\rule{\linewidth}{0.6pt}\par}\vspace{0.6em}}

% ---- 記號巨集（共用的；其餘在專案的 .sty）----
\newcommand{\dd}{\mathop{}\!\mathrm{d}}
\newcommand{\e}{\mathrm{e}}
\newcommand{\avg}[1]{\left\langle #1\right\rangle}
\newcommand{\Order}{\mathcal O}

% ---- 區塊共用計數器：定義、定理、推導、例題共用 (3.k) ----
\newcounter{qblock}
\renewcommand{\theqblock}{\lecturechapter.\arabic{qblock}}
\newcommand{\qtitle}[1]{\ifstrempty{#1}{}{\enspace#1}}

% ---- 不讓標題落在頁底 ----
\makeatletter
\newcommand{\qblockspace}[1]{\if@nobreak\else\Needspace*{#1}\fi}
\makeatother
\newcommand{\coursekeepblocks}{%
  \pretocmd{\section}{\qblockspace{12\baselineskip}}{}{}%
  \pretocmd{\subsection}{\qblockspace{8\baselineskip}}{}{}%
}

% ---- 區塊 ----
% 框內 \label 記到區塊編號：tcolorbox 在 \sbox 裡做 \refstepcounter，\@currentlabel 出了那個群組就還原，
% 框內的 \label 因此記成外層的小節號（DE L02 遇過，2026-10-21）。\qsetlabel 在框的正文開頭重設一次；
% 自己寫 before upper 的編號框要把它放進去。沒有編號的框不受影響。
\makeatletter
\newcommand{\qsetlabel}{\ifdefined\tcbcounter\ifx\tcbcounter\@empty\else
  \protected@edef\@currentlabel{\lecturechapter.\thetcbcounter}\fi\fi}
\makeatother
\tcbset{qbase/.style={enhanced,breakable,arc=2pt,boxrule=0.6pt,
  left=8pt,right=8pt,top=5pt,bottom=5pt,before skip=0.9em,after skip=0.9em,
  fonttitle=\bfseries\small,toptitle=2pt,bottomtitle=2pt,
  before upper={\setlength{\parskip}{0.45em}\qsetlabel}}}
\tcbset{qbar/.style={enhanced,breakable,frame hidden,sharp corners,
  left=10pt,right=8pt,top=5pt,bottom=5pt,before skip=0.8em,after skip=0.8em,
  before upper={\setlength{\parskip}{0.45em}\qsetlabel}}}
\newtcolorbox[use counter=qblock]{definition}[1][]{qbase,
  colframe=qblue,colback=qblue!4,coltitle=white,colbacktitle=qblue,
  title={定義 \lecturechapter.\thetcbcounter\qtitle{#1}}}
\newtcolorbox[use counter=qblock]{theorem}[1][]{qbase,
  colframe=qpurple,colback=qpurple!4,coltitle=white,colbacktitle=qpurple,
  title={定理 \lecturechapter.\thetcbcounter\qtitle{#1}}}
\newtcolorbox{note}[1][]{qbase,
  colframe=qorange,colback=qorangebg,coltitle=white,colbacktitle=qorange,
  title={注意\qtitle{#1}}}
\newtcolorbox[use counter=qblock]{derivation}[1][]{qbar,colback=white,top=1pt,bottom=1pt,
  borderline west={1.2pt}{0pt}{qrule},
  before upper={\setlength{\parskip}{0.45em}\qsetlabel{\small\bfseries\color{qnavy}推導 \lecturechapter.\thetcbcounter\qtitle{#1}}\par}}
\newtcolorbox[use counter=qblock]{example}[1][]{qbar,colback=white,top=1pt,bottom=1pt,
  borderline west={1.2pt}{0pt}{qgreen!60},
  before upper={\setlength{\parskip}{0.45em}\qsetlabel{\small\bfseries\color{qgreen}例 \lecturechapter.\thetcbcounter\qtitle{#1}}\par}}
\newtcolorbox{extension}[1][]{qbar,colback=qgraybg,fontupper=\small\color{qmuted},
  before upper={{\bfseries 延伸\qtitle{#1}}\par}}
% 重點結果：與詳解相同的淺橘底＋左側橘條
\newtcolorbox{qresult}{enhanced,frame hidden,sharp corners,colback=qorangebg,
  borderline west={2.2pt}{0pt}{qorange},left=10pt,right=10pt,top=4pt,bottom=4pt,
  before skip=0.7em,after skip=0.7em,halign=center}
\newcommand{\keyresult}[1]{\par\begin{qresult}\makebox[\linewidth]{$\displaystyle #1$}\end{qresult}}
% 章首問題：淺灰底
\newcommand{\chaptergoal}[1]{%
  \begin{tcolorbox}[qbar,colback=qgraybg]{\small\bfseries\color{qnavy}本章問題}\par #1\end{tcolorbox}}

% 例題的解：接在 example 之後，綠色小標，不加框
\newenvironment{solution}
  {\par\smallskip\noindent{\small\bfseries\color{qgreen}解}\enspace\ignorespaces}
  {\par\medskip}
% 例題、推導前留空間，避免小標落在頁底
\BeforeBeginEnvironment{example}{\qblockspace{6\baselineskip}}
\BeforeBeginEnvironment{derivation}{\qblockspace{6\baselineskip}}
\BeforeBeginEnvironment{definition}{\qblockspace{5\baselineskip}}
\BeforeBeginEnvironment{theorem}{\qblockspace{5\baselineskip}}
% 可跨頁的框剛好在頁底結束時，xdvipdfmx 的顏色堆疊可能錯位，之後的正文變成白字
% （Neamen L05_v1 遇過，2026-10-15）。框的環境群組結束後重設一次正文顏色；
% 放在框的 after 鍵裡沒有用，因為群組結束時顏色又被恢復成錯位的狀態。
\AfterEndEnvironment{definition}{\normalcolor}
\AfterEndEnvironment{theorem}{\normalcolor}
\AfterEndEnvironment{property}{\normalcolor}
\AfterEndEnvironment{note}{\normalcolor}
\AfterEndEnvironment{derivation}{\normalcolor}
\AfterEndEnvironment{example}{\normalcolor}
\AfterEndEnvironment{extension}{\normalcolor}
}{\input{"D:/Codex workplace/_templates/styles/codex_notes.tex"}}
% 參數（都可省略）：\codexsecsep（節號分隔，預設 -）、\codexlecturetag{章號}（頁首右側，預設 L03）、\codextoclabelwidth（預設 2.6em）。
% 版面：A4、10pt（主檔 \documentclass[10pt,a4paper]{article}）、四邊 2cm、行距 1.3、Noto Sans TC、段落不縮排；深藍標題、藍色節號、淺橘重點結果。
% 巨集：\setuplecture{章號}{章名}、\lecturehead{小字}{章名}、\booksection{中文}{English}、\plainsection、\chaptergoal、\keyresult（只放數學式）、
% definition／theorem／note／derivation／example／solution／extension[標題]、\coursekeepblocks；記號 \dd、\e、\avg、\Order。
\RequirePackage[no-math]{fontspec}
\RequirePackage{xeCJK}
% 英文撇號：' 經 TeX 連字變成 U+2019，xeCJK 預設把 ‘’ 當全形標點，Noto Sans TC 的這兩個字形也是全寬，
% Ampère's 的撇號後面因此多一段空隙。把 ‘’ 改回西文字元，主字型加 pwid 取比例寬字形（只作用於西文字型，
% 中文標點由 CJK 字型排，不受影響）。中文引號用「」，“” 維持原狀。2026-10-22
\defaultfontfeatures+[\rmfamily]{RawFeature=+pwid}
\xeCJKDeclareCharClass{Default}{"2018,"2019}
% 字型資料夾：載入本檔前 \def\codexfontpath{<資料夾>/} 可改；設成空白則依檔名在 TeX 搜尋路徑與目前資料夾找。2026-11-15
\providecommand{\codexfontpath}{D:/Codex workplace/Noto_Sans_TC/static/}
\setmainfont{NotoSansTC-Regular.ttf}[
  Path=\codexfontpath,
  BoldFont=NotoSansTC-Bold.ttf,
  ItalicFont=NotoSansTC-Regular.ttf,
  BoldItalicFont=NotoSansTC-Bold.ttf
]
\setCJKmainfont{NotoSansTC-Regular.ttf}[
  Path=\codexfontpath,
  BoldFont=NotoSansTC-Bold.ttf,
  ItalicFont=NotoSansTC-Regular.ttf,
  BoldItalicFont=NotoSansTC-Bold.ttf
]
\setmonofont{NotoSansTC-Regular.ttf}[
  Path=\codexfontpath,
  BoldFont=NotoSansTC-Bold.ttf
]
\setCJKmonofont{NotoSansTC-Regular.ttf}[
  Path=\codexfontpath,
  BoldFont=NotoSansTC-Bold.ttf
]
\RequirePackage{amsmath,amssymb,amsthm,mathtools}
\RequirePackage[a4paper,margin=2cm]{geometry}
\RequirePackage{booktabs,longtable,array,tabularx}
\RequirePackage{graphicx,float}
\RequirePackage{enumitem,setspace,xcolor,fancyhdr,needspace,etoolbox}
\RequirePackage[most]{tcolorbox}
\RequirePackage{titlesec,titletoc}
\RequirePackage{caption}
\RequirePackage{tikz,pgfplots}
\RequirePackage[unicode]{hyperref}
\pgfplotsset{compat=1.18}
\usetikzlibrary{arrows.meta,calc,positioning,decorations.pathreplacing}
\usepgfplotslibrary{fillbetween}

% ---- 色票（與共用詳解版型相同，另加綠、紫各一）----
\definecolor{qblue}{HTML}{1D5FA8}
\definecolor{qnavy}{HTML}{173D6B}
\definecolor{qorange}{HTML}{EB6834}
\definecolor{qorangebg}{HTML}{FFF4E8}
\definecolor{qgraybg}{HTML}{F3F4F6}
\definecolor{qrule}{HTML}{D8DDE3}
\definecolor{qmuted}{HTML}{6B7280}
\definecolor{qtext}{HTML}{1F2329}
\definecolor{qgreen}{HTML}{2E7D4F}
\definecolor{qpurple}{HTML}{6B3FA0}
% 舊版型用到的灰色名稱（圖檔可能引用）
\colorlet{qgray}{qmuted}
\colorlet{qlight}{qgraybg}
\hypersetup{colorlinks=true,linkcolor=qnavy,urlcolor=qblue,citecolor=qblue}

% ---- 段落 ----
\setstretch{1.3}
\setlength{\parindent}{0pt}
\setlength{\parskip}{0.55em}
\setlength{\headheight}{16pt}
\setlength{\headsep}{0.45cm}
\setlength{\footskip}{1cm}
\setlist{itemsep=0.2em,topsep=0.3em}
\setlist[itemize]{label=\textcolor{qblue}{\small\textbullet}}
\setlist[enumerate]{label=\textcolor{qblue}{\arabic*.}}
\renewcommand{\arraystretch}{1.16}
\allowdisplaybreaks[2]
\raggedbottom
\widowpenalty=10000
\clubpenalty=10000
\emergencystretch=3em
\AtBeginDocument{\color{qtext}}

% ---- 圖說 ----
\renewcommand{\figurename}{圖}
\renewcommand{\tablename}{表}
\captionsetup{font={small,color=qmuted},labelfont={bf,color=qnavy},labelsep=quad,
  justification=justified,singlelinecheck=true,skip=6pt}

% ---- 章設定：節號 3-1（分隔符 \codexsecsep）、式號與區塊號 (3.k)；頁首右側 \codexlecturetag{章號}　章名 ----
\providecommand{\codexsecsep}{-}
\providecommand{\codexlecturetag}[1]{L\ifnum#1<10 0\fi#1}
\providecommand{\codextoclabelwidth}{2.6em}
\newcommand{\lecturechapter}{0}
\newcommand{\lecturetitle}{}
\newcommand{\setuplecture}[2]{%
  \renewcommand{\lecturechapter}{#1}%
  \renewcommand{\lecturetitle}{#2}%
  \renewcommand{\thesection}{#1\codexsecsep\arabic{section}}%
  \renewcommand{\theequation}{#1.\arabic{equation}}%
  \renewcommand{\thefigure}{#1.\arabic{figure}}%
  \renewcommand{\thetable}{#1.\arabic{table}}%
  \fancyhead[L]{\color{qmuted}\small\coursename}%
  \fancyhead[R]{\color{qmuted}\small \codexlecturetag{#1}\quad #2}%
}

% ---- 頁首頁尾 ----
\pagestyle{fancy}
\fancyhf{}
\fancyfoot[C]{\color{qmuted}\small\thepage}
\renewcommand{\headrulewidth}{0pt}

% ---- 標題 ----
\titleformat{\section}[block]{\color{qnavy}\Large\bfseries\boldmath}
  {\textcolor{qblue}{\thesection}}{0.8em}{}
\titleformat{name=\section,numberless}[block]{\color{qnavy}\Large\bfseries\boldmath}{}{0pt}{}
\titlespacing*{\section}{0pt}{2.2em}{0.8em}
\titleformat{\subsection}{\color{qnavy}\normalsize\bfseries\boldmath}{}{0pt}{}
\titlespacing*{\subsection}{0pt}{1.2em}{0.25em}

% 課本節：\booksection{中文標題}{English title}；英文副標放同一行，放不下整段換行，不斷字
\newsavebox{\qsubtitlebox}
\DeclareRobustCommand{\qbooksubtitle}[1]{%
  \sbox{\qsubtitlebox}{\normalfont\small(#1)}%
  \ifdim\wd\qsubtitlebox<\dimexpr\linewidth-4em\relax
    \hskip 0pt plus 1fil\penalty0\hskip 0.8em plus -1fil\relax
    {\normalfont\small\color{qmuted}#1}%
  \else
    \enspace{\normalfont\small\color{qmuted}
      \ifcsname l@nohyphenation\endcsname\language\csname l@nohyphenation\endcsname\fi
      #1}%
  \fi}
\newcommand{\booksection}[2]{\section[#1]{#1\qbooksubtitle{#2}}}
\newcommand{\plainsection}[1]{\section*{#1}\addcontentsline{toc}{section}{#1}}

% ---- 目錄 ----
\renewcommand{\contentsname}{目錄}
\setcounter{tocdepth}{1}
\titlecontents{section}[\codextoclabelwidth]{\addvspace{0.35em}\color{qtext}}
  {\contentslabel[\textcolor{qblue}{\bfseries\thecontentslabel}]{\codextoclabelwidth}}{}
  {\color{qmuted}\titlerule*[0.6pc]{.}\contentspage}

% ---- 章首標頭（取代 \maketitle）----
\newcommand{\lecturehead}[2]{%
  {\color{qmuted}\small #1\par}\vspace{2pt}
  {\color{qnavy}\fontsize{18pt}{24pt}\selectfont\bfseries #2\par}\vspace{6pt}
  {\color{qrule}\rule{\linewidth}{0.6pt}\par}\vspace{0.6em}}

% ---- 記號巨集（共用的；其餘在專案的 .sty）----
\newcommand{\dd}{\mathop{}\!\mathrm{d}}
\newcommand{\e}{\mathrm{e}}
\newcommand{\avg}[1]{\left\langle #1\right\rangle}
\newcommand{\Order}{\mathcal O}

% ---- 區塊共用計數器：定義、定理、推導、例題共用 (3.k) ----
\newcounter{qblock}
\renewcommand{\theqblock}{\lecturechapter.\arabic{qblock}}
\newcommand{\qtitle}[1]{\ifstrempty{#1}{}{\enspace#1}}

% ---- 不讓標題落在頁底 ----
\makeatletter
\newcommand{\qblockspace}[1]{\if@nobreak\else\Needspace*{#1}\fi}
\makeatother
\newcommand{\coursekeepblocks}{%
  \pretocmd{\section}{\qblockspace{12\baselineskip}}{}{}%
  \pretocmd{\subsection}{\qblockspace{8\baselineskip}}{}{}%
}

% ---- 區塊 ----
% 框內 \label 記到區塊編號：tcolorbox 在 \sbox 裡做 \refstepcounter，\@currentlabel 出了那個群組就還原，
% 框內的 \label 因此記成外層的小節號（DE L02 遇過，2026-10-21）。\qsetlabel 在框的正文開頭重設一次；
% 自己寫 before upper 的編號框要把它放進去。沒有編號的框不受影響。
\makeatletter
\newcommand{\qsetlabel}{\ifdefined\tcbcounter\ifx\tcbcounter\@empty\else
  \protected@edef\@currentlabel{\lecturechapter.\thetcbcounter}\fi\fi}
\makeatother
\tcbset{qbase/.style={enhanced,breakable,arc=2pt,boxrule=0.6pt,
  left=8pt,right=8pt,top=5pt,bottom=5pt,before skip=0.9em,after skip=0.9em,
  fonttitle=\bfseries\small,toptitle=2pt,bottomtitle=2pt,
  before upper={\setlength{\parskip}{0.45em}\qsetlabel}}}
\tcbset{qbar/.style={enhanced,breakable,frame hidden,sharp corners,
  left=10pt,right=8pt,top=5pt,bottom=5pt,before skip=0.8em,after skip=0.8em,
  before upper={\setlength{\parskip}{0.45em}\qsetlabel}}}
\newtcolorbox[use counter=qblock]{definition}[1][]{qbase,
  colframe=qblue,colback=qblue!4,coltitle=white,colbacktitle=qblue,
  title={定義 \lecturechapter.\thetcbcounter\qtitle{#1}}}
\newtcolorbox[use counter=qblock]{theorem}[1][]{qbase,
  colframe=qpurple,colback=qpurple!4,coltitle=white,colbacktitle=qpurple,
  title={定理 \lecturechapter.\thetcbcounter\qtitle{#1}}}
\newtcolorbox{note}[1][]{qbase,
  colframe=qorange,colback=qorangebg,coltitle=white,colbacktitle=qorange,
  title={注意\qtitle{#1}}}
\newtcolorbox[use counter=qblock]{derivation}[1][]{qbar,colback=white,top=1pt,bottom=1pt,
  borderline west={1.2pt}{0pt}{qrule},
  before upper={\setlength{\parskip}{0.45em}\qsetlabel{\small\bfseries\color{qnavy}推導 \lecturechapter.\thetcbcounter\qtitle{#1}}\par}}
\newtcolorbox[use counter=qblock]{example}[1][]{qbar,colback=white,top=1pt,bottom=1pt,
  borderline west={1.2pt}{0pt}{qgreen!60},
  before upper={\setlength{\parskip}{0.45em}\qsetlabel{\small\bfseries\color{qgreen}例 \lecturechapter.\thetcbcounter\qtitle{#1}}\par}}
\newtcolorbox{extension}[1][]{qbar,colback=qgraybg,fontupper=\small\color{qmuted},
  before upper={{\bfseries 延伸\qtitle{#1}}\par}}
% 重點結果：與詳解相同的淺橘底＋左側橘條
\newtcolorbox{qresult}{enhanced,frame hidden,sharp corners,colback=qorangebg,
  borderline west={2.2pt}{0pt}{qorange},left=10pt,right=10pt,top=4pt,bottom=4pt,
  before skip=0.7em,after skip=0.7em,halign=center}
\newcommand{\keyresult}[1]{\par\begin{qresult}\makebox[\linewidth]{$\displaystyle #1$}\end{qresult}}
% 章首問題：淺灰底
\newcommand{\chaptergoal}[1]{%
  \begin{tcolorbox}[qbar,colback=qgraybg]{\small\bfseries\color{qnavy}本章問題}\par #1\end{tcolorbox}}

% 例題的解：接在 example 之後，綠色小標，不加框
\newenvironment{solution}
  {\par\smallskip\noindent{\small\bfseries\color{qgreen}解}\enspace\ignorespaces}
  {\par\medskip}
% 例題、推導前留空間，避免小標落在頁底
\BeforeBeginEnvironment{example}{\qblockspace{6\baselineskip}}
\BeforeBeginEnvironment{derivation}{\qblockspace{6\baselineskip}}
\BeforeBeginEnvironment{definition}{\qblockspace{5\baselineskip}}
\BeforeBeginEnvironment{theorem}{\qblockspace{5\baselineskip}}
% 可跨頁的框剛好在頁底結束時，xdvipdfmx 的顏色堆疊可能錯位，之後的正文變成白字
% （Neamen L05_v1 遇過，2026-10-15）。框的環境群組結束後重設一次正文顏色；
% 放在框的 after 鍵裡沒有用，因為群組結束時顏色又被恢復成錯位的狀態。
\AfterEndEnvironment{definition}{\normalcolor}
\AfterEndEnvironment{theorem}{\normalcolor}
\AfterEndEnvironment{property}{\normalcolor}
\AfterEndEnvironment{note}{\normalcolor}
\AfterEndEnvironment{derivation}{\normalcolor}
\AfterEndEnvironment{example}{\normalcolor}
\AfterEndEnvironment{extension}{\normalcolor}
}{\input{"D:/Codex workplace/_templates/styles/codex_notes.tex"}}
% 參數（都可省略）：\codexsecsep（節號分隔，預設 -）、\codexlecturetag{章號}（頁首右側，預設 L03）、\codextoclabelwidth（預設 2.6em）。
% 版面：A4、10pt（主檔 \documentclass[10pt,a4paper]{article}）、四邊 2cm、行距 1.3、Noto Sans TC、段落不縮排；深藍標題、藍色節號、淺橘重點結果。
% 巨集：\setuplecture{章號}{章名}、\lecturehead{小字}{章名}、\booksection{中文}{English}、\plainsection、\chaptergoal、\keyresult（只放數學式）、
% definition／theorem／note／derivation／example／solution／extension[標題]、\coursekeepblocks；記號 \dd、\e、\avg、\Order。
\RequirePackage[no-math]{fontspec}
\RequirePackage{xeCJK}
% 英文撇號：' 經 TeX 連字變成 U+2019，xeCJK 預設把 ‘’ 當全形標點，Noto Sans TC 的這兩個字形也是全寬，
% Ampère's 的撇號後面因此多一段空隙。把 ‘’ 改回西文字元，主字型加 pwid 取比例寬字形（只作用於西文字型，
% 中文標點由 CJK 字型排，不受影響）。中文引號用「」，“” 維持原狀。2026-10-22
\defaultfontfeatures+[\rmfamily]{RawFeature=+pwid}
\xeCJKDeclareCharClass{Default}{"2018,"2019}
% 字型資料夾：載入本檔前 \def\codexfontpath{<資料夾>/} 可改；設成空白則依檔名在 TeX 搜尋路徑與目前資料夾找。2026-11-15
\providecommand{\codexfontpath}{D:/Codex workplace/Noto_Sans_TC/static/}
\setmainfont{NotoSansTC-Regular.ttf}[
  Path=\codexfontpath,
  BoldFont=NotoSansTC-Bold.ttf,
  ItalicFont=NotoSansTC-Regular.ttf,
  BoldItalicFont=NotoSansTC-Bold.ttf
]
\setCJKmainfont{NotoSansTC-Regular.ttf}[
  Path=\codexfontpath,
  BoldFont=NotoSansTC-Bold.ttf,
  ItalicFont=NotoSansTC-Regular.ttf,
  BoldItalicFont=NotoSansTC-Bold.ttf
]
\setmonofont{NotoSansTC-Regular.ttf}[
  Path=\codexfontpath,
  BoldFont=NotoSansTC-Bold.ttf
]
\setCJKmonofont{NotoSansTC-Regular.ttf}[
  Path=\codexfontpath,
  BoldFont=NotoSansTC-Bold.ttf
]
\RequirePackage{amsmath,amssymb,amsthm,mathtools}
\RequirePackage[a4paper,margin=2cm]{geometry}
\RequirePackage{booktabs,longtable,array,tabularx}
\RequirePackage{graphicx,float}
\RequirePackage{enumitem,setspace,xcolor,fancyhdr,needspace,etoolbox}
\RequirePackage[most]{tcolorbox}
\RequirePackage{titlesec,titletoc}
\RequirePackage{caption}
\RequirePackage{tikz,pgfplots}
\RequirePackage[unicode]{hyperref}
\pgfplotsset{compat=1.18}
\usetikzlibrary{arrows.meta,calc,positioning,decorations.pathreplacing}
\usepgfplotslibrary{fillbetween}

% ---- 色票（與共用詳解版型相同，另加綠、紫各一）----
\definecolor{qblue}{HTML}{1D5FA8}
\definecolor{qnavy}{HTML}{173D6B}
\definecolor{qorange}{HTML}{EB6834}
\definecolor{qorangebg}{HTML}{FFF4E8}
\definecolor{qgraybg}{HTML}{F3F4F6}
\definecolor{qrule}{HTML}{D8DDE3}
\definecolor{qmuted}{HTML}{6B7280}
\definecolor{qtext}{HTML}{1F2329}
\definecolor{qgreen}{HTML}{2E7D4F}
\definecolor{qpurple}{HTML}{6B3FA0}
% 舊版型用到的灰色名稱（圖檔可能引用）
\colorlet{qgray}{qmuted}
\colorlet{qlight}{qgraybg}
\hypersetup{colorlinks=true,linkcolor=qnavy,urlcolor=qblue,citecolor=qblue}

% ---- 段落 ----
\setstretch{1.3}
\setlength{\parindent}{0pt}
\setlength{\parskip}{0.55em}
\setlength{\headheight}{16pt}
\setlength{\headsep}{0.45cm}
\setlength{\footskip}{1cm}
\setlist{itemsep=0.2em,topsep=0.3em}
\setlist[itemize]{label=\textcolor{qblue}{\small\textbullet}}
\setlist[enumerate]{label=\textcolor{qblue}{\arabic*.}}
\renewcommand{\arraystretch}{1.16}
\allowdisplaybreaks[2]
\raggedbottom
\widowpenalty=10000
\clubpenalty=10000
\emergencystretch=3em
\AtBeginDocument{\color{qtext}}

% ---- 圖說 ----
\renewcommand{\figurename}{圖}
\renewcommand{\tablename}{表}
\captionsetup{font={small,color=qmuted},labelfont={bf,color=qnavy},labelsep=quad,
  justification=justified,singlelinecheck=true,skip=6pt}

% ---- 章設定：節號 3-1（分隔符 \codexsecsep）、式號與區塊號 (3.k)；頁首右側 \codexlecturetag{章號}　章名 ----
\providecommand{\codexsecsep}{-}
\providecommand{\codexlecturetag}[1]{L\ifnum#1<10 0\fi#1}
\providecommand{\codextoclabelwidth}{2.6em}
\newcommand{\lecturechapter}{0}
\newcommand{\lecturetitle}{}
\newcommand{\setuplecture}[2]{%
  \renewcommand{\lecturechapter}{#1}%
  \renewcommand{\lecturetitle}{#2}%
  \renewcommand{\thesection}{#1\codexsecsep\arabic{section}}%
  \renewcommand{\theequation}{#1.\arabic{equation}}%
  \renewcommand{\thefigure}{#1.\arabic{figure}}%
  \renewcommand{\thetable}{#1.\arabic{table}}%
  \fancyhead[L]{\color{qmuted}\small\coursename}%
  \fancyhead[R]{\color{qmuted}\small \codexlecturetag{#1}\quad #2}%
}

% ---- 頁首頁尾 ----
\pagestyle{fancy}
\fancyhf{}
\fancyfoot[C]{\color{qmuted}\small\thepage}
\renewcommand{\headrulewidth}{0pt}

% ---- 標題 ----
\titleformat{\section}[block]{\color{qnavy}\Large\bfseries\boldmath}
  {\textcolor{qblue}{\thesection}}{0.8em}{}
\titleformat{name=\section,numberless}[block]{\color{qnavy}\Large\bfseries\boldmath}{}{0pt}{}
\titlespacing*{\section}{0pt}{2.2em}{0.8em}
\titleformat{\subsection}{\color{qnavy}\normalsize\bfseries\boldmath}{}{0pt}{}
\titlespacing*{\subsection}{0pt}{1.2em}{0.25em}

% 課本節：\booksection{中文標題}{English title}；英文副標放同一行，放不下整段換行，不斷字
\newsavebox{\qsubtitlebox}
\DeclareRobustCommand{\qbooksubtitle}[1]{%
  \sbox{\qsubtitlebox}{\normalfont\small(#1)}%
  \ifdim\wd\qsubtitlebox<\dimexpr\linewidth-4em\relax
    \hskip 0pt plus 1fil\penalty0\hskip 0.8em plus -1fil\relax
    {\normalfont\small\color{qmuted}#1}%
  \else
    \enspace{\normalfont\small\color{qmuted}
      \ifcsname l@nohyphenation\endcsname\language\csname l@nohyphenation\endcsname\fi
      #1}%
  \fi}
\newcommand{\booksection}[2]{\section[#1]{#1\qbooksubtitle{#2}}}
\newcommand{\plainsection}[1]{\section*{#1}\addcontentsline{toc}{section}{#1}}

% ---- 目錄 ----
\renewcommand{\contentsname}{目錄}
\setcounter{tocdepth}{1}
\titlecontents{section}[\codextoclabelwidth]{\addvspace{0.35em}\color{qtext}}
  {\contentslabel[\textcolor{qblue}{\bfseries\thecontentslabel}]{\codextoclabelwidth}}{}
  {\color{qmuted}\titlerule*[0.6pc]{.}\contentspage}

% ---- 章首標頭（取代 \maketitle）----
\newcommand{\lecturehead}[2]{%
  {\color{qmuted}\small #1\par}\vspace{2pt}
  {\color{qnavy}\fontsize{18pt}{24pt}\selectfont\bfseries #2\par}\vspace{6pt}
  {\color{qrule}\rule{\linewidth}{0.6pt}\par}\vspace{0.6em}}

% ---- 記號巨集（共用的；其餘在專案的 .sty）----
\newcommand{\dd}{\mathop{}\!\mathrm{d}}
\newcommand{\e}{\mathrm{e}}
\newcommand{\avg}[1]{\left\langle #1\right\rangle}
\newcommand{\Order}{\mathcal O}

% ---- 區塊共用計數器：定義、定理、推導、例題共用 (3.k) ----
\newcounter{qblock}
\renewcommand{\theqblock}{\lecturechapter.\arabic{qblock}}
\newcommand{\qtitle}[1]{\ifstrempty{#1}{}{\enspace#1}}

% ---- 不讓標題落在頁底 ----
\makeatletter
\newcommand{\qblockspace}[1]{\if@nobreak\else\Needspace*{#1}\fi}
\makeatother
\newcommand{\coursekeepblocks}{%
  \pretocmd{\section}{\qblockspace{12\baselineskip}}{}{}%
  \pretocmd{\subsection}{\qblockspace{8\baselineskip}}{}{}%
}

% ---- 區塊 ----
% 框內 \label 記到區塊編號：tcolorbox 在 \sbox 裡做 \refstepcounter，\@currentlabel 出了那個群組就還原，
% 框內的 \label 因此記成外層的小節號（DE L02 遇過，2026-10-21）。\qsetlabel 在框的正文開頭重設一次；
% 自己寫 before upper 的編號框要把它放進去。沒有編號的框不受影響。
\makeatletter
\newcommand{\qsetlabel}{\ifdefined\tcbcounter\ifx\tcbcounter\@empty\else
  \protected@edef\@currentlabel{\lecturechapter.\thetcbcounter}\fi\fi}
\makeatother
\tcbset{qbase/.style={enhanced,breakable,arc=2pt,boxrule=0.6pt,
  left=8pt,right=8pt,top=5pt,bottom=5pt,before skip=0.9em,after skip=0.9em,
  fonttitle=\bfseries\small,toptitle=2pt,bottomtitle=2pt,
  before upper={\setlength{\parskip}{0.45em}\qsetlabel}}}
\tcbset{qbar/.style={enhanced,breakable,frame hidden,sharp corners,
  left=10pt,right=8pt,top=5pt,bottom=5pt,before skip=0.8em,after skip=0.8em,
  before upper={\setlength{\parskip}{0.45em}\qsetlabel}}}
\newtcolorbox[use counter=qblock]{definition}[1][]{qbase,
  colframe=qblue,colback=qblue!4,coltitle=white,colbacktitle=qblue,
  title={定義 \lecturechapter.\thetcbcounter\qtitle{#1}}}
\newtcolorbox[use counter=qblock]{theorem}[1][]{qbase,
  colframe=qpurple,colback=qpurple!4,coltitle=white,colbacktitle=qpurple,
  title={定理 \lecturechapter.\thetcbcounter\qtitle{#1}}}
\newtcolorbox{note}[1][]{qbase,
  colframe=qorange,colback=qorangebg,coltitle=white,colbacktitle=qorange,
  title={注意\qtitle{#1}}}
\newtcolorbox[use counter=qblock]{derivation}[1][]{qbar,colback=white,top=1pt,bottom=1pt,
  borderline west={1.2pt}{0pt}{qrule},
  before upper={\setlength{\parskip}{0.45em}\qsetlabel{\small\bfseries\color{qnavy}推導 \lecturechapter.\thetcbcounter\qtitle{#1}}\par}}
\newtcolorbox[use counter=qblock]{example}[1][]{qbar,colback=white,top=1pt,bottom=1pt,
  borderline west={1.2pt}{0pt}{qgreen!60},
  before upper={\setlength{\parskip}{0.45em}\qsetlabel{\small\bfseries\color{qgreen}例 \lecturechapter.\thetcbcounter\qtitle{#1}}\par}}
\newtcolorbox{extension}[1][]{qbar,colback=qgraybg,fontupper=\small\color{qmuted},
  before upper={{\bfseries 延伸\qtitle{#1}}\par}}
% 重點結果：與詳解相同的淺橘底＋左側橘條
\newtcolorbox{qresult}{enhanced,frame hidden,sharp corners,colback=qorangebg,
  borderline west={2.2pt}{0pt}{qorange},left=10pt,right=10pt,top=4pt,bottom=4pt,
  before skip=0.7em,after skip=0.7em,halign=center}
\newcommand{\keyresult}[1]{\par\begin{qresult}\makebox[\linewidth]{$\displaystyle #1$}\end{qresult}}
% 章首問題：淺灰底
\newcommand{\chaptergoal}[1]{%
  \begin{tcolorbox}[qbar,colback=qgraybg]{\small\bfseries\color{qnavy}本章問題}\par #1\end{tcolorbox}}

% 例題的解：接在 example 之後，綠色小標，不加框
\newenvironment{solution}
  {\par\smallskip\noindent{\small\bfseries\color{qgreen}解}\enspace\ignorespaces}
  {\par\medskip}
% 例題、推導前留空間，避免小標落在頁底
\BeforeBeginEnvironment{example}{\qblockspace{6\baselineskip}}
\BeforeBeginEnvironment{derivation}{\qblockspace{6\baselineskip}}
\BeforeBeginEnvironment{definition}{\qblockspace{5\baselineskip}}
\BeforeBeginEnvironment{theorem}{\qblockspace{5\baselineskip}}
% 可跨頁的框剛好在頁底結束時，xdvipdfmx 的顏色堆疊可能錯位，之後的正文變成白字
% （Neamen L05_v1 遇過，2026-10-15）。框的環境群組結束後重設一次正文顏色；
% 放在框的 after 鍵裡沒有用，因為群組結束時顏色又被恢復成錯位的狀態。
\AfterEndEnvironment{definition}{\normalcolor}
\AfterEndEnvironment{theorem}{\normalcolor}
\AfterEndEnvironment{property}{\normalcolor}
\AfterEndEnvironment{note}{\normalcolor}
\AfterEndEnvironment{derivation}{\normalcolor}
\AfterEndEnvironment{example}{\normalcolor}
\AfterEndEnvironment{extension}{\normalcolor}
